%&pdflatex
\documentclass[a4paper]{article}

\usepackage[english,spanish]{babel}

\usepackage{amsmath,amsfonts,amsthm}

\newtheorem{theorem}{Teorema}
\newtheorem*{corollary*}{Corolario}
\newtheorem{lemma}{Lema}

\newcommand{\N}{\mathbb{N}}
\newcommand{\calO}{\mathcal{O}}

\def\limtoenorme{\mathop{\,\relbar\joinrel\relbar\joinrel%
  \longrightarrow\,}\limits}

\date{}

%%%%%%%%%%%%%%%%%%%%%%%%%%%%%%%%%%%%%%%%%%%%%%%%%%%%%%%%%%%%%%%
\title{ESTIMACIONES DE POLINOMIOS ORTOGONALES CON PESO
$(1-x)^{\alpha} (1+x)^{\beta} + M\delta_{-1} + N\delta_{1}$}

\author{Mario P\'erez Riera y Juan Luis Varona Malumbres%
\renewcommand\thefootnote{}
\renewcommand\footnotemark{}
\footnote{
Este art\'{\i}culo se present\'o en las
\textit{XIII Jornadas Hispano Lusas de Matem\'aticas}, 
Universidad de Valladolid (Valladolid, 1988).
Contrariamente a lo anunciado ~---y dado el tiempo ya transcurrido---
parece claro que no va a haber nunca actas de dicho congreso.
As\'{\i}, \'este es un manuscrito sin publicar; no obstante,
el resultado principal de este trabajo ~---el teorema~2---
tambi\'en est\'a demostrado, con un m\'etodo totalmente distinto
y mucho m\'as general al aqu\'{\i} empleado,
en el art\'{\i}culo
\textsc{J.~J. Guadalupe, M.~P\'erez, F.~J.~Ruiz y J.~L.~Varona},
{Asymptotic behaviour of orthogonal polynomials relative to
measures with mass points},
\textit{Mathematika}
\textbf{40} (1993), 331--344.
}
}

%%%%%%%%%%%%%%%%%%%%%%%%%%%%%%%%%%%%%%%%%%%%%%%%%%%%%%%%%%%%%%%

%%%%%%%%%%%%%%%%%%%%%%%%%%%%%%%%%%%%%%%%%%%%%%%%%%%%%%%%%%%%%%%
\begin{document}
%%%%%%%%%%%%%%%%%%%%%%%%%%%%%%%%%%%%%%%%%%%%%%%%%%%%%%%%%%%%%%%

%%%%%%%%%%%%%%%%%%%%%%%%%%%%%%%%%%%%%%%%%%%%%%%%%%%%%%%%%%%%%%%
\maketitle
%%%%%%%%%%%%%%%%%%%%%%%%%%%%%%%%%%%%%%%%%%%%%%%%%%%%%%%%%%%%%%%


%%%%%%%%%%%%%%%%%%%%%%%%%%%%%%%%%%%%%%%%%%%%%%%%%%%%%%%%%%%%%%%
\selectlanguage{english}
\begin{abstract}\noindent
Let $w$ be a Jacobi weight function and $\delta_{t}$ the 
delta function at a point~$t$. We obtain some estimates for the 
polynomials which are orthogonal with respect to $w + M \delta_{-1} 
+ N \delta_{1}$, by using expresions which relate them to the 
Jacobi polynomials.
\end{abstract}
\selectlanguage{spanish}
%%%%%%%%%%%%%%%%%%%%%%%%%%%%%%%%%%%%%%%%%%%%%%%%%%%%%%%%%%%%%%%

%%%%%%%%%%%%%%%%%%%%%%%%%%%%%%%%%%%%%%%%%%%%%%%%%%%%%%%%%%%%%%%
\begin{quote}
\textit{Clasificaci\'on A.M.S. (1980)}: 42C10.
\end{quote}
\bigskip
%%%%%%%%%%%%%%%%%%%%%%%%%%%%%%%%%%%%%%%%%%%%%%%%%%%%%%%%%%%%%%%

Sean $\alpha, \beta > -1$ y $\{P^{(\alpha,\beta)}_{n}\}_{n \in \N}$
los polinomios de Jacobi usuales~\cite{Sz}, ortogonales con respecto 
a la funci\'{o}n peso $(1-x)^{\alpha}(1+x)^{\beta}$ sobre el 
intervalo $[-1,1]$. Algunas propiedades que necesitaremos son
Comencemos con
$$
P^{(\alpha,\beta)}_{n}(1) = \frac{(\alpha+1)_{n}}{n!} ,
\quad\text{donde}\quad
(\alpha+1)^{n} = \frac{\Gamma(\alpha+n+1)}{\Gamma(\alpha+1)},
$$
$$
P^{(\alpha,\beta)}_{n}(-x) = (-1)^{n}P^{(\beta,\alpha)}_{n}(x).
$$
El coeficiente director del polinomio $P^{(\alpha,\beta)}_{n}$ es
$$
k^{(\alpha,\beta)}_{n} = \frac{2^{-n} \Gamma(2n+\alpha+\beta+1)}{\Gamma(n+1)\Gamma(n+\alpha+\beta+1)},
$$
y su norma al cuadrado vale
\begin{align*}
h^{(\alpha,\beta)}_{n} &= \int_{-1}^{1} [P^{(\alpha,\beta)}_{n}(x)]^{2} 
(1-x)^{\alpha}(1+x)^{\beta} \,dx \\
&= \frac{2^{\alpha+\beta+1}\Gamma(n+\alpha+1)\Gamma(n+\beta+1)}
  {(2n+\alpha+\beta+1)\Gamma(n+1)\Gamma(n+\alpha+\beta+1)}.
\end{align*}
F\'{o}rmula de Rodrigues:
$$
P^{(\alpha,\beta)}_{n}(x) 
= (-2)^{-n}(n!)^{-1} (1-x)^{-\alpha}(1+x)^{-\beta} \, \frac{d^{n}}{dx^{n}} 
[(1-x)^{n+\alpha}(1+x)^{n+\beta}].
$$
Finalmente, es conocida~\cite{Mu} la acotaci\'{o}n
$$
(h^{(\alpha,\beta)}_{n})^{-1/2} |P^{(\alpha,\beta)}_{n}(x)| 
\leq K \left(1-x + \frac{1}{n^{2}} \right)^{-\alpha/2-1/4}
\left(1+x + \frac{1}{n^{2}} \right)^{-\beta/2-1/4}.
$$

Sean ahora $M,N \geq 0$ y
$$
d\mu(x) = \frac{\Gamma(\alpha+\beta+2)}
  {2^{\alpha+\beta+1}\Gamma(\alpha+1)\Gamma(\beta+1)}
(1-x)^{\alpha} (1+x)^{\beta} \,dx + M \delta_{-1}(x) + N \delta_{1}(x)
$$
sobre el intervalo $[-1,1]$, donde $\delta_{\pm1}$ 
es la delta de Dirac sobre el punto~$\pm1$. El prop\'{o}sito 
de este trabajo es obtener las mismas acotaciones para los polinomios 
ortonormales con respecto a~$d\mu$.

Sea
\begin{equation}
\label{eq:1}
\begin{aligned}
&P^{(\alpha,\beta,M,N)}_{n}(x) =
\left[\frac{(\alpha+\beta+1)_{n}}{n!}\right]^{2} 
\bigg\{ \frac{n+\alpha+\beta+1}{2(\alpha+\beta+1)} \\
&\qquad\null\times [B_{n}M(1-x)-A_{n}N(1+x)] P^{(\alpha+1,\beta+1)}_{n-1}(x) 
+ A_{n}B_{n}P^{(\alpha,\beta)}_{n}(x) \bigg\},
\end{aligned}
\end{equation}
donde
$$
A_{n} = \frac{(\alpha+1)_{n} n!}{(\beta+1)_{n}(\alpha+\beta+1)_{n}}
+ \frac{n(n+\alpha+\beta+1)M}{(\beta+1)(\alpha+\beta+1)}
$$
y
$$
B_{n} = \frac{(\beta+1)_{n} n!}{(\alpha+1)_{n}(\alpha+\beta+1)_{n}} 
+ \frac{n(n+\alpha+\beta+1)N}{(\alpha+1)(\alpha+\beta+1)}
$$
(el caso $\alpha+\beta+1 = 0$ debe ser entendido por continuidad 
en $\alpha$ y~$\beta$). Koornwinder~\cite{Ko} demuestra que estos 
polinomios son ortogonales en $[-1,1]$ con respecto a la medida~$d\mu$.

Llamaremos $k^{(\alpha,\beta,M,N)}_{n}$ al 
coeficiente director de $P^{(\alpha,\beta,M,N)}_{n}$ 
y $h^{(\alpha,\beta,M,N)}_{n}$ al cuadrado 
de su norma:
$$
h^{(\alpha,\beta,M,N)}_{n}
= \int_{-1}^{1} [P^{(\alpha,\beta,M,N)}_{n}(x)]^{2}\,d\mu(x).
$$
Usaremos asimismo la notaci\'{o}n $r_{n} \sim s_{n}$ 
cuando $r_{n}/s_{n} \limtoenorme_{n\to\infty} 1$,
siendo $\{r_{n}\}$ y $\{s_{n}\}$ sucesiones de t\'{e}rminos no nulos. 
Mediante la f\'{o}rmula de Stirling puede probarse, a partir de~\eqref{eq:1}:
\begin{equation}
k^{(\alpha,\beta,M,N)}_{n} \sim 
\begin{cases}
2^{n+\alpha+\beta} \pi^{-1/2} n^{-1/2} &\text{si } M=N=0, \\[3pt]
\frac{N \Gamma(\beta+1)}{\pi^{1/2}\Gamma(\alpha+\beta+2)\Gamma(\alpha+2)} 
\, 2^{n+\alpha+\beta} n^{2+2\alpha-1/2} &\text{si } N>M=0, \\[4pt]
\frac{M \Gamma(\alpha+1)}{\pi^{1/2}\Gamma(\alpha+\beta+2)\Gamma(\beta+2)} 
\, 2^{n+\alpha+\beta} n^{2+2\beta-1/2} &\text{si } M>N=0, \\[3pt]
\frac{M N}{\pi^{1/2}(\alpha+1)(\beta+1)\Gamma(\alpha+\beta+2)^{2}} 
\, 2^{n+\alpha+\beta} n^{4+2\alpha+2\beta-1/2} &\text{si } M,N>0, 
\end{cases}
\label{eq:2}
\end{equation}
\begin{equation}
P^{(\alpha,\beta,M,N)}_{n}(1) \sim 
\frac{1}{\Gamma(\alpha+1)} \, n^{\alpha}
+ \frac{M}{\Gamma(\beta+2)\Gamma(\alpha+\beta+2)} \, n^{2\beta+\alpha+2},
\label{eq:3}
\end{equation}
\begin{equation}
(-1)^{n}P^{(\alpha,\beta,M,N)}_{n}(-1) \sim 
\frac{1}{\Gamma(\beta+1)} \, n^{\beta}
+ \frac{N}{\Gamma(\alpha+2)\Gamma(\alpha+\beta+2)} \, n^{2\alpha+\beta+2}.
\label{eq:4}
\end{equation}
Para el c\'{a}lculo de $h^{(\alpha,\beta,M,N)}_{n}$ 
necesitamos algunos resultados previos:

\begin{lemma}
\label{lem:1}
$$
\int_{-1}^{1} P^{(\alpha+1,\beta+1)}_{n-1}(x) (1-x)^{\alpha+1}(1+x)^{\beta} \,dx 
\sim (-1)^{n-1} 2^{\alpha+\beta+2} \Gamma(\beta+1) n^{-(\beta+1)}.
$$
\end{lemma}

\begin{proof}
Se deduce integrando por partes, despu\'{e}s 
de sustituir la f\'{o}rmula de Rodrigues.
\end{proof}

\begin{lemma}
\label{lem:2}
\begin{multline*}
\int_{-1}^{1} P^{(\alpha+1,\beta+1)}_{n-1}(x)
[(1-x)^{n}-2^{n}](1-x)^{\alpha+1}(1+x)^{\beta} \,dx \\
= -\int_{-1}^{1} P^{(\alpha+1,\beta+1)}_{n-1}(x)
(1-x)^{n-1}(1-x)^{\alpha+1}(1+x)^{\beta} \,dx 
\sim \frac{(-1)^{n}\pi^{1/2}}{2^{n-1}} \, n^{-1/2}.
\end{multline*}
\end{lemma}

\begin{proof}
Se sigue de la igualdad
$(1-x)^{n}-2^{n} = -(1+x) \sum_{k=0}^{n-1} 2^{k}(1-x)^{n-1-k}$.
\end{proof}

Utilizando \eqref{eq:1}, \eqref{eq:2}, \eqref{eq:3}, \eqref{eq:4}
y los lemas~\ref{lem:1} y~\ref{lem:2}, se puede ahora 
demostrar:

\begin{theorem}
\label{th:1}
$$
h^{(\alpha,\beta,M,N)}_{n} \sim
\begin{cases}
\frac{\Gamma(\alpha+\beta+2)}{2\Gamma(\alpha+1)\Gamma(\beta+1)} 
\, n^{-1} & \text{si } M=N=0, \\[3pt]
\frac{N^{2} \Gamma(\beta+1)}{2\Gamma(\alpha+1)\Gamma(\alpha+2)^{2} \Gamma(\alpha+\beta+2)} 
\, n^{3+4\alpha} & \text{si } N>M=0, \\[3pt]
\frac{M^{2} \Gamma(\alpha+1)}{2\Gamma(\beta+1)\Gamma(\beta+2)^{2} \Gamma(\alpha+\beta+2)} 
\, n^{3+4\beta} & \text{si } M>N=0, \\[3pt]
\frac{M^{2}N^{2}}{2(\alpha+1)(\beta+1)\Gamma(\alpha+2)\Gamma(\beta+2) \Gamma(\alpha+\beta+2)^{3}} 
\, n^{7+4\alpha+4\beta} & \text{si } M,N>0.
\end{cases}
$$
\end{theorem}

\begin{corollary*}
Si llamamos $p_{n}(x)$ a los polinomios ortonormales 
con respecto a~$d\mu$, entonces
$$
p_{n}(1) \sim
\begin{cases}
\frac{2^{1/2}\Gamma(\beta+1)^{1/2}}{\Gamma(\alpha+1)^{1/2} \Gamma(\alpha+\beta+2)^{1/2}} 
\, n^{\alpha+1/2} & \text{si } N=0, \\[4pt]
\frac{2^{1/2}(\alpha+1)\Gamma(\alpha+1)^{1/2} \Gamma(\alpha+\beta+2)^{1/2}}
{N \Gamma(\beta+1)^{1/2}} \, n^{-\alpha-3/2} & \text{si } N>0.
\end{cases}
$$
Un resultado an\'{a}logo se obtiene para $(-1)^{n}p_{n}(-1)$.
\end{corollary*}

Finalmente, obtenemos para los polinomios $P^{(\alpha,\beta,M,N)}_{n}$ 
la cota que busc\'{a}bamos:

\begin{theorem}
Sean $\alpha, \beta > -1$, $M,N \geq 0$.
Existe una constante $K$ tal que $\forall x \in [-1,1]$, $\forall n \geq 0$:
\begin{multline*}
(h^{(\alpha,\beta,M,N)}_{n})^{-1/2} |P^{(\alpha,\beta,M,N)}_{n}(x)| \\
\leq K \left(1-x + \frac{1}{n^{2}} \right)^{-\alpha/2-1/4}
\left(1+x + \frac{1}{n^{2}} \right)^{-\beta/2-1/4}.
\end{multline*}
\end{theorem}

\begin{proof}
Por simetr\'{\i}a, basta demostrar que $\forall x \geq 0$
el miembro de la izquierda est\'{a} acotado por
$K (1-x + {1}/{n^{2}} )^{-\alpha/2-1/4}$. Como en el caso de 
los polinomios de Jacobi~\cite{Mu}, ello ser\'{a} cierto si se tiene
$$
(h^{(\alpha,\beta,M,N)}_{n})^{-1/2} |P^{(\alpha,\beta,M,N)}_{n}(\cos \theta)|
=
\begin{cases}
\calO(\theta)^{-\alpha-1/2}) & \text{si } c/n \leq \theta \leq \pi/2, \\
\calO(n^{\alpha+1/2}) & \text{si } 0 \leq \theta \leq c/n, 
\end{cases}
$$
donde $c$ es una constante positiva prefijada cualquiera. Como 
esta expresi\'{o}n es v\'{a}lida para los polinomios de
Jacobi~\cite[T.~7.32.2]{Sz}, el problema puede reducirse,
por~\eqref{eq:1}, a la acotaci\'{o}n de
\begin{multline*}
(h^{(\alpha,\beta,M,N)}_{n})^{-1/2} 
\left[\frac{(\alpha+\beta+1)_{n}}{n!}\right]^{2}
A_{n} \bigg\{B_{n}P^{(\alpha,\beta)}_{n}(x) \\
- N \frac{n+\alpha+\beta+1}{2(\alpha+\beta+1)} (1+x) 
P^{(\alpha+1,\beta+1)}_{n-1}(x) \bigg\}.
\end{multline*}
Utilizando ahora el teorema~\ref{th:1}, estimaciones para $A_{n}$
y $(\alpha+\beta+1)_{n}/{n!}$, las igualdades
$$
(1+x) P^{(\alpha+1,\beta+1)}_{n-1}(x) 
= \frac{2(n+\beta)}{2n+\alpha+\beta+1}  P^{(\alpha+1,\beta)}_{n-1}(x) 
+ \frac{2n}{2n+\alpha+\beta+1}  P^{(\alpha+1,\beta)}_{n}(x) 
$$
(\cite[4.5.4]{Sz}), y
$$
P^{(\alpha,\beta)}_{n}(x) = \frac{n+\alpha+\beta+1}{2n+\alpha+\beta+1} 
P^{(\alpha+1,\beta)}_{n}(x) - \frac{n+\beta}{2n+\alpha+\beta+1} 
P^{(\alpha+1,\beta)}_{n-1}(x) ,
$$
y la f\'{o}rmula de tipo Hilb \cite[T.~8.21.12]{Sz}
\begin{multline*}
\left(\sen \frac{\theta}{2}\right)^{\alpha+1}
\left(\cos \frac{\theta}{2}\right)^{\beta} 
 P^{(\alpha+1,\beta)}_{n}(\cos \theta) = \\
D^{-(\alpha+1)} \frac{\Gamma(n+\alpha+2)}{n!}
\left(\frac{\theta}{\sen \theta}\right)^{1/2} J_{\alpha+1}(D\theta) \\
+
\begin{cases}
\theta^{1/2} \calO(n^{-3/2}) & \text{si } c/n \leq \theta \leq \pi/2, \\
\theta^{\alpha+3}\calO(n^{\alpha+1}) & \text{si } 0 \leq \theta \leq c/n,
\end{cases}
\end{multline*}
donde $D = n + (\alpha+\beta+2)/2$ y $J_{\alpha+1}(z)$ 
es la funci\'{o}n de Bessel, puede verse que el teorema est\'{a} 
probado si se demuestra que las dos expresiones
\begin{equation}
\label{eq:5}
\begin{aligned}
&D^{-2\alpha-2} \, \frac{(n+\alpha+1)(n+\alpha+\beta+1)}{n} 
\bigg[ \frac{(\beta+1)_{n} n!}{(\alpha+1)_{n}(\alpha+\beta+1)_{n}} \\
&\quad\null+ N \frac{n(n+\beta)}{(\alpha+1)(\alpha+\beta+1)} \bigg]
\left| D^{\alpha+1} J_{\alpha+1}(D\theta) - 
(D-1)^{\alpha+1} J_{\alpha+1}((D-1)\theta)) \right|
\end{aligned}
\end{equation}
y
\begin{equation}
\label{eq:6}
\begin{aligned}
&D^{-2\alpha-2} \Bigg| \frac{(n+\alpha+{1})(n+\alpha+\beta+1)}{n}
(D-1)^{\alpha+1} \\
&\qquad\qquad\null \times
\bigg[ \frac{(\beta+1)_{n} n!}{(\alpha+1)_{n}(\alpha+\beta+1)_{n}} 
+ N \frac{n(n+\beta)}{(\alpha+1)(\alpha+\beta+1)} \bigg] \\
&\qquad\qquad\qquad\null- (n+\beta)\frac{D^{2\alpha+2}}{(D-1)^{\alpha+1}} 
\bigg[ \frac{(\beta+1)_{n}n!}{(\alpha+1)_{n}(\alpha+\beta+1)_{n}} \\
&\qquad\qquad\qquad\qquad\null+ N \frac{(n+\alpha+1)(n+\alpha+\beta+1)}
{(\alpha+1)(\alpha+\beta+1)} \bigg]
\Bigg| J_{\alpha+1}((D-1)\theta)
\end{aligned}
\end{equation}
son ambas
$$
n^{3/2-\alpha} \theta^{\alpha+1}
\begin{cases}
\calO(\theta^{-\alpha-1/2}) & \text{si } c/n \leq \theta \leq \pi/2, \\
\calO(n^{\alpha+1/2}) & \text{si } 0 \leq \theta \leq c/n.
\end{cases}
$$
Para lograr estas estimaciones, empleamos algunas propiedades 
de las funciones de Bessel:
\begin{gather*}
[z^{\nu}J_{\nu}(z)]' = z^{\nu}J_{\nu-1}(z); \\
J_{\nu}(z) = \calO(z^{\nu}) \quad\text{cuando}\quad z \to 0;
\qquad
J_{\nu}(z) = \calO(z^{-1/2}) \quad\text{cuando}\quad z \to \infty.
\end{gather*}
De estas propiedades y del teorema del valor medio, puede deducirse 
que
\begin{multline*}
\left| (D\theta)^{\alpha+1} J_{\alpha+1}(D\theta) - ((D-1)\theta)^{\alpha+1}
J_{\alpha+1}((D-1)\theta)) \right| \\
= n^{\alpha+1/2} \theta^{2\alpha+2}
\begin{cases}
\calO(\theta^{-\alpha-1/2}) & \text{si } c/n \leq \theta \leq \pi/2, \\
\calO(n^{\alpha+1/2}) & \text{si } 0 \leq \theta \leq c/n,
\end{cases}
\end{multline*}
y, de aqu\'{\i}, la estimaci\'{o}n asint\'{o}tica~\eqref{eq:5}. Por \'{u}ltimo, 
\eqref{eq:6} se logra mediante las expresiones asint\'{o}ticas citadas de 
$J_{\nu}$ y utilizando la f\'{o}rmula
\begin{multline*}
\frac{(\beta+1)_{n} n!}{(\alpha+1)_{n}(\alpha+\beta+1)_{n}} \\
= \frac{\Gamma(\alpha+1)\Gamma(\alpha+\beta+1)}{\Gamma(\beta+1)} 
\, n^{-2\alpha} \left[1 - \frac{\alpha(\alpha+\beta+1)}{n} 
+ \calO(n^{-2}) \right],
\end{multline*}
que puede deducirse de la expresi\'{o}n~\cite{Er}:
$$
\frac{\Gamma(n+a)}{\Gamma(n+b)} 
= n^{a-b} \left[1 + \frac{(a-b)(a+b-1)}{2n} + \calO(n^{-2}) \right].
$$
\end{proof}

%%%%%%%%%%%%%%%%%%%%%%%%%%%%%%%%%%%%%%%%%%%%%%%%%%%%%%%%%%%%%%%

\renewcommand{\refname}{Bibliograf\'{\i}a}

\frenchspacing

\begin{thebibliography}{4}

\frenchspacing

\bibitem{Er}
%[1]\tab 
\textsc{A. Erd\'{e}lyi y otros},
``Higher transcendental functions'', 
McGraw-Hill Book Co., New York, vol.~I, 1953.

\bibitem{Ko}
%[2]\tab 
\textsc{T. H. Koornwinder},
Orthogonal polynomials with weight function 
$(1-x)^{\alpha} (1+x)^{\beta} + M \delta(x+1) + N \delta(x-1)$,
\textit{Canad. Math. Bull.} \textbf{27} (1984), 205--214.

\bibitem{Mu}
%[3]\tab 
\textsc{B. Muckenhoupt}, Mean convergence of Jacobi series,
\textit{Proc. Amer. Math. Soc.} \textbf{23} (1969), 306--310.

\bibitem{Sz}
%[4]\tab 
\textsc{G. Szeg\H{o}}, ``Orthogonal Polynomials'', 3.\textsuperscript{a}~ed.,
Amer. Math. Soc. Colloq. Publ., vol. 23, Amer. Math. Soc., 
Providence, R.I., 1967.

\end{thebibliography}

%%%%%%%%%%%%%%%%%%%%%%%%%%%%%%%%%%%%%%%%%%%%%%%%%%%%%%%%%%%%%%%

\bigskip
\noindent
\begin{tabular}{@{}l}
Mario P\'{e}rez Riera\\
Dpto. de Matem\'{a}ticas\\
Facultad de Ciencias\\
Univ. de Zaragoza\\
50009 Zaragoza (Espa\~{n}a)
\end{tabular}\hfill
\begin{tabular}{l@{}}
Juan Luis Varona Malumbres\\
Dpto. de Matem\'{a}ticas\\
Col. Univ. de La Rioja\\
Univ. de Zaragoza\\
26001 Logro\~{n}o (La Rioja, Espa\~{n}a)
\end{tabular}


\end{document}
